% problem_id: S014-lurie-ultracategories-corollary-8-3-5
% source_id: S014 (Jacob Lurie, Ultracategories)
% source_file: lurie-conceptual.pdf
% statement_locator: printed p. 105
% proof_locator: printed p. [105]
% representation: PDF; transcribed from rendered page images (never from the text layer)
% status: complete (statement + author proof verbatim + reason + locators)
%
%% ============================================================================
%% Verbatim transcription (read off rendered page images, NOT off the text layer)
%% ----------------------------------------------------------------------------
%% Source paper : J. Lurie, "Ultracategories"
%% Host path    : /disks/sata1/yupeng/human-proof-corpus/source-cache/registry-sources/lurie/lurie-conceptual.pdf
%% PDF facts    : 147 pages, 612x792 pt (letter), pdfTeX-1.40.18, created 2020-02-06
%% Target item  : Corollary 8.3.5
%% Statement    : printed page 105 = PDF page 105 (page header reads "ULTRACATEGORIES 105")
%% Proof        : printed page 105 = PDF page 105 (same page; proof closes with □ on that page)
%% ----------------------------------------------------------------------------
%% Locating step (pdftotext ONLY; its text layer mangles symbols, e.g. "∈" -> "2",
%% "α" -> "↵", so it was used for positions/line numbers only, never as source text):
%%   pdftotext <pdf> /tmp/lurie.txt
%%   "Corollary 8.3.5." occurs at line 6099, "Proof of Corollary 8.3.5." at line 6108,
%%   both located on PDF page 105; the forward-reference
%%   "ultrastructures, which we will prove in §8.3 (see Corollary 8.3.5)." is at line 1673 (Remark 2.4.6).
%% Rendering step (the images actually read):
%%   pdftoppm -f 104 -l 106 -r 150 -png <pdf> /tmp/m1/M1
%%   pdftoppm -f 105 -l 105 -r 300 -png -x 200 -y 250 -W 2200 -H 700 <pdf> /tmp/m1/M1-105-crop
%%   pdftoppm -f 105 -l 105 -r 300 -png -x 200 -y 900 -W 2200 -H 700 <pdf> /tmp/m1/M1-105-proof
%%   (copies kept in ./renders/)
%% ----------------------------------------------------------------------------
%% Conventions used below:
%%  * The wording is transcribed literally: nothing rewritten, simplified or completed.
%%  * The theorem/proof environments are transcriber-supplied containers only; the
%%    running text inside them is verbatim. The original sets the corollary body in
%%    italic and the proof in upright roman with an italic "Proof of ..." head and a
%%    trailing □, which is what these containers reproduce.
%%  * Suspect printing in the original is copied as printed and flagged with a comment.
%%  * Everything printed in the original that belongs to the target is on page 105.
%% ============================================================================

\documentclass[11pt]{article}
\usepackage{amsmath,amssymb,amsthm}

\newcommand{\Ult}{\mathrm{Ult}}
\newcommand{\Fun}{\mathrm{Fun}}

\theoremstyle{plain}
\newtheorem*{corollaryM}{Corollary 8.3.5}
\theoremstyle{definition}
\newtheorem*{remarkM}{Remark 8.3.6}

\begin{document}

%% --- page 105 ---

\begin{corollaryM}
Let $\mathcal{M}$ and $\mathcal{N}$ be categories which admit small products and small
filtered colimits, and egard $\mathcal{M}$ and $\mathcal{N}$ as equipped with the categorical
ultrastructures of Example 1.3.8. %% [PRINTED AS IN ORIGINAL: "egard" — sic, the "r" is missing; read "regard". Copied verbatim, not silently corrected.]
Let $\Fun^{\Ult,\Pi}(\mathcal{M},\mathcal{N})$ denote the full subcategory of
$\Fun^{\Ult}(\mathcal{M},\mathcal{N})$ spanned by those functors which preserve small products.
Then the forgetful functor $\Fun^{\Ult,\Pi}(\mathcal{M},\mathcal{N}) \to \Fun(\mathcal{M},\mathcal{N})$
is a fully faithful embedding, whose essential image is spanned by those ultrafunctors
$F : \mathcal{M} \to \mathcal{N}$ which preserve small products and small filtered colimits.
\end{corollaryM}

\begin{proof}[Proof of Corollary 8.3.5]
It follows from Corollary 8.3.2 that the forgetful functor
$\theta : \Fun^{\Ult,\Pi}(\mathcal{M},\mathcal{N}) \to \Fun(\mathcal{M},\mathcal{N})$
is fully faithful. Any functor $F : \mathcal{M} \to \mathcal{N}$ belonging to the essential image
of $\theta$ must preserve small products (by definition) and small filtered colimits
(by Proposition 5.3.4). Conversely, if $F : \mathcal{M} \to \mathcal{N}$ preserves small products
and small filtered colimits, then it admits an ultrastructure by virtue of Proposition 1.4.9,
and therefore belongs to the essential image of $\theta$.
\end{proof}

%% ============================================================================
%% CONTEXT BELOW — NOT part of the target item; printed on the same page 105,
%% immediately after the proof. Included because it is the informal restatement
%% of the deferred claim that Corollary 8.3.5 discharges, and it is on a page
%% that was image-verified.
%% ============================================================================

\begin{remarkM}
Corollary 8.3.5 implies in particular that if $F : \mathcal{M} \to \mathcal{N}$ is a functor
which preserves small products and small filtered colimits, then it admits a \textit{unique}
ultrastructure (namely, the ultrastructure given by Proposition 1.4.9).
\end{remarkM}

%% ============================================================================
%% LOCATING-TRAIL ONLY — text-layer derived pointers, NOT transcribed from images,
%% NOT to be quoted as the original wording. Recorded so the deferral chain is auditable.
%%
%%  * Remark 2.4.6 (text layer, line 1673; PDF page 29 by form-feed count — the page
%%    image was NOT read, so the page number is unverified): ends with
%%    "... This is a special case of a more general result about categorical
%%    ultrastructures, which we will prove in §8.3 (see Corollary 8.3.5)."
%%    => this is the deferral whose target is the item transcribed above.
%%  * Remark 1.3.8 (text layer, line 926): contains the pointer
%%    "... the categorical ultrastructure is "initial" among all possible
%%    ultrastructures on M (see Example 8.3.4)."
%%  * §8.3 opening (text layer, line 6015-6018, printed page 103): "... In this section,
%%    we prove a more general form of this result (Corollary 8.3.5), which we deduce from
%%    a more general statement in the setting of right ultrastructures."
%%    "deduce from" is the author's own description of the proof strategy — and the
%%    transcribed proof bears this out: it is a three-sentence deduction by citation.
%%  * Proof of Corollary 8.3.2 (printed page 104, image-verified as the page preceding the
%%    target; NOT transcribed here) is where the substantive argument sits: it reduces to
%%    Theorem 8.2.5 and Proposition 8.3.1 (replacing theta by a restriction functor), then
%%    proves (i) via a right Kan extension and (ii') via an isomorphism criterion.
%% ============================================================================

\end{document}
